% STUFF
\clearpage
\section{Outer Solver}

%We consider a resource allocation problem where the goal is to optimise a black-box objective function \( f: \mathcal{X} \to \mathbb{R} \) using Bayesian Optimisation.
%The function \( f \) is expensive to evaluate and only accessible via noisy observations. 
%The domain \( \mathcal{X} \subseteq \mathbb{R}^d \) represents a constrained space defined by hardware inventory limits.
%Our objective is to find an optimal configuration of GPU allocations that maximizes the function \( f \).

%
%\subsubsection*{Bayesian Optimization Formulation}
%
%Bayesian Optimization models the objective function \( f \) using a probabilistic surrogate model \( \hat{f}(x) \), typically a Gaussian Process (GP):
%
%\[
%\hat{f}(x) \sim \mathcal{GP}(\mu(x), k(x, x'))
%\]
%
%An acquisition function \( \alpha(x; \hat{f}) \) is used to determine the next point to evaluate:
%
%\[
%x_{t+1} = \arg \max_{x \in \mathcal{X}} \alpha(x; \hat{f}_t)
%\]
%
%This optimization is subject to a set of resource constraints, which are defined below.

%We refer to the bounded bin packing problem as the \textit{inner loop}, and the outer optimisation process as the \textit{outer loop}, as illustrated in \autoref{fig:04-problem-formulation}.
%The \textit{outer loop} is responsible for generating promising island layouts, which are then passed to the \textit{inner loop} for resource packing.
%The resulting throughput is returned to the \textit{outer loop}, allowing the optimisation process to evaluate and refine its search for effective configurations.
% This iterative cycle can be repeated multiple times to generate progressively generate an optimum solution.
 
% During \autoref{chap:constraint-optimisation} we formulated a non-linear bounded bin-packing problem, by constraining the islands to a pre-defined type and size.
%In a true scheduling problem we expect to be allocated an inventory of different GPU types, not formed islands.
%Let \(T\) be the set of GPU types an for each \(t\in T\), let \(n_t\) be the number of GPUs of type \(t\) allocated in the inventory.
%The full GPU inventory is then the vector \(\bigl(n_t\bigr)_{t\in T}\), with total number of GPUs in the system modelled as \(N = \sum_{t\in T} n_t\).

%The \textit{inner loop} receives an island layout proposal from the \textit{outer loop}, this layout is specified as a finite index set \(\mathcal{I} \;=\; \{\,i = 1,\dots,|\mathcal{I}|\}\), where each island \(i\) is represented by a pair \(\bigl(t_{i},s_{i}\bigr)\).  
%Here, \(t_{i}\in T\) denotes the GPU type assigned to island \(i\), and \(s_{i}\in \mathbb{Z}_{\ge 0}\) indicates the number of GPUs of type \(t_{i}\) allocated to that island.
%The \textit{outer loop} guarantees that these allocations respect the total-inventory constraints.

%Given the predefined island layout \(\bigl\{(t_{j},s_{j})\bigr\}_{j=1}^{|\mathcal{I}|}\), the inner loop's only task is to assign workloads of different prefill and decode lengths.
%No further modification of \(t_{i}\) or \(s_{i}\) takes place at this stage, as the inner loop focuses entirely on determining the optimal bin-packing configuration.

The \textit{outer loop} receives an input inventory, containing a list of GPU types (\(t\in T\)) and quantities (\(n_t\)), represented with the set \(\{\,I_t : t \in T\}\).
Given the inventory, the \textit{outer loop's} job is to propose new island layouts (\(\mathcal{I}\)), which can be subsequently evaluated by the \textit{inner loop}.

The \textit{inner loop} execution is somewhat expensive, hence a complete grid search would incur significant compute time.
Instead we use a refined Bayesian Optimisation search method with two different methods for searching, \textit{direct} (\autoref{chap:island-direct} and \textit{divided} (\autoref{chap:island-divide}).

\begin{figure}[ht]
    \centering
    \includegraphics[width=\textwidth,height=5cm]{example-image-a}
    \caption[Diagram showing effective purpose of \textit{outer loop}]
    {Diagram showing effective purpose of \textit{outer loop}.}
    \label{fig:04-outer-loop-demo}
\end{figure}

We select Bayesian Optimisation specifically due to its ability to efficiently explore high-dimensional, expensive‐to‐evaluate spaces with only a few hundred evaluations, rather than the exponential cost of a full grid search.
By maintaining a surrogate model of the island‐layout performance (Gaussian process over layouts), Bayesian Optimisation balances exploration (sampling layouts in poorly understood regions) with exploitation (refining around the current best layouts).

Unlike reinforcement learning, which often requires thousands of interactions to learn a policy, BO is designed for “one‐shot” black‐box optimisation.
With the correct parameters it can deliver near‐optimal solutions with far fewer simulator calls, minimal hyper-parameter tuning, and strong convergence.

% First BO Method
\subsection{Direct Island Generation}
\label{chap:island-direct}

In the \textit{direct} island generation approach, islands are explicitly defined within the Bayesian Optimisation process.
Configurations correspond to a fixed number of islands, with each candidate solution representing a specific allocation strategy.
The number of GPUs in each island is directly part of the output space, where specific variables represent the GPU count per island `slot'.
This setup provides precise control over how GPUs are grouped, but it also requires well-defined constraints to ensure that all generated configurations are valid.

\subsubsection*{Search Parameters}

\begin{itemize}
    \item \( T \): the set of GPU types, e.g., \( T = \{\text{A100}, \text{V100}, \ldots\} \)
    \item \( K \): the number of GPU slots per type
    \item \( x_{t,k} \): the integer decision variable representing the number of GPUs of type \( t \in T \) assigned to slot \( k \in \{1, \ldots, K\} \)
    \item \( I_t \): the total inventory of GPU type \( t \in T \)
\end{itemize}

%Each parameter \( x_{k,t} \) is bounded:
%
%\[
%x_{g,k} \in \{0, 1, \ldots, I_t\}, \quad \forall t \in T, \; \forall k \in \{1, \ldots, K\}
%\]

\subsubsection*{Search Constraints}

For each GPU type \( t \in T \), we enforce that the total number of GPUs allocated across all \( K \) slots does not exceed the inventory \( I_t \):
\[
\sum_{k=1}^{K} x_{g,k} \leq I_t, \quad \forall t \in T
\]

\subsubsection*{Implementation}

\todo{
\begin{itemize}
\item Ax Implementation.
\item Why SAASBO may be required (high dimensionality)
\item Why this approach is challenging.
\end{itemize}
}

% Though BO has been recognized as a powerful approaches for a diversity of problem domains, it is also known that traditional BO is inherently vulnerable to the ”curse-of- dimensionality” [6] when the search space has a large number of dimensions. wl-446

%These constraints define a feasible subspace \( \mathcal{X} \subseteq \mathbb{R}^d \), over which the Bayesian Optimization is performed.

%\subsubsection*{Search Space Definition}
%
%In implementation (e.g., using Ax or BoTorch), the search space \( \mathcal{X} \) is constructed by defining:
%\begin{itemize}
%    \item A list of bounded integer parameters \( \{x_{g,k}\} \)
%    \item A set of linear \texttt{SumConstraint} objects that enforce the inventory limits
%\end{itemize}
%
%These constraints are handled internally by the BO engine when sampling candidate configurations during optimization.


% Second BO method
\clearpage
\subsection{Divider Island Generation}
\label{chap:island-divide}

The \textit{divider} island generation method uses a more abstract representation that substantially reduces the dimensionality of the BO problem.
Explicit allocation at the slot level requires one parameter per slot per GPU type, resulting in a parameter space of size \( |G| \times K \).
This approach becomes quickly computationally expensive due to the high number of search parameters and the need for constraints to ensure valid configurations.

To address this, the divider-based approach reduces the search space to just two parameters per GPU type: the total number of islands and a scalar that controls the balance of GPU allocation across those islands.
 This formulation results in a significantly smaller parameter space of size \( |G| \times 2 \), enabling more efficient optimisation and faster convergence.

An advantage of this method is that it avoids the need for explicit resource constraints, which are otherwise necessary to enforce inventory limits and prevent invalid configurations.
By shifting from low-level allocation to higher-level distribution strategies, divider island generation maintains control over performance-influencing characteristics while offering a more scalable and interpretable optimisation process.

%The primary goal of the island divider procedure is to programmatically divide a given total number of GPUs of a specific type into a list of smaller, integer-sized islands.

\subsubsection{Input Variables}
Let \(G_{\text{total}}\) denote the total number of GPUs available for that type, and let \(N_{\text{target}}\) represent the desired number of islands to create.
We can allow the actual number of islands may be lower if there are not enough GPUs to satisfy the minimum-size requirement.
The parameter \(S_{\text{skew}}\) controls how any leftover GPUs (after each island has been assigned its minimum) are distributed: when \(S_{\text{skew}}\approx0.0\), the extra GPUs are spread as evenly as possible; if \(S_{\text{skew}}>0.0\), islands with higher indices receive proportionally more of the remainder; and if \(S_{\text{skew}}<0.0\), islands with lower indices receive more.
Finally, \(I_{\text{min}}\) specifies the minimum number of GPUs that any island must contain.

% we repeat this for every GPU type

\begin{figure}[ht]
    \centering
    \includegraphics[width=\textwidth]{divider_skew_effects_overlay_islands.pdf}
    \caption[Graph displaying the effect of changing the \(S_{\text{skew}}\) value]
    {Graph displaying the effect of changing the \(S_{\text{skew}}\) value where \(G_{\text{total}}=100\) and \(N_{\text{target}}=5\).}
    \label{fig:04-skew-changes}
\end{figure}

\subsubsection{Preconditions}
Initially, we verify that the input parameters meet a series of constraints, if any of these conditions is not met, the procedure raises an error and terminates immediately.
We require \(I_{\text{min}} > 0\), \(G_{\text{total}}\) and \(N_{\text{target}}\) must be nonnegative.
Next, we perform trivial checks to assess whether it is possible to create at least one island given the input parameters.
If \(N_{\text{target}} = 0\), \(G_{\text{total}} = 0\) or \(G_{\text{total}} < I_{\text{min}}\), then no islands are to be created and the function returns the empty set \(\varnothing\).
Similarly, if , there are insufficient GPUs to form even a single island of size \(I_{\text{min}}\), so the function again returns \(\varnothing\).

\subsubsection{Number of Islands}

To establish how many islands can feasibly be created given the available GPUs and the minimum-size requirement, we first compute \(N_{\text{max\_possible}} = \lfloor G_{\text{total}} / I_{\text{min}} \rfloor\), which represents the maximum number of islands if each island is allocated exactly \(I_{\text{min}}\) GPUs.
We then set \(N_{\text{actual}} = \min\bigl(N_{\text{target}},\,N_{\text{max\_possible}}\bigr)\), so that the actual number of islands cannot exceed either the user's target or the physical limit imposed by the GPU pool.
If \(N_{\text{actual}} = 0\), this indicates that even a single island of size \(I_{\text{min}}\) cannot be formed (either because \(N_{\text{target}} = 0\) or because \(G_{\text{total}} < I_{\text{min}}\)), and the function returns the empty set \(\varnothing\).
Otherwise, \(N_{\text{actual}}\) islands will be created, each initially allocated \(I_{\text{min}}\) GPUs before any further distribution of the remaining GPUs occurs.

\subsubsection{Base Allocation}

In the base allocation step, we allocate \(I_{\text{min}}\) GPUs to each of the \(N_{\text{actual}}\) islands, so that the total GPUs used for this initial distribution is \(G_{\text{base\_used}} = N_{\text{actual}} \cdot I_{\text{min}}\).
The number of GPUs remaining to be distributed is then \(G_{\text{remaining}} = G_{\text{total}} - G_{\text{base\_used}}\).
At this point, every island has been provisioned with its minimum allocation of \(I_{\text{min}}\) GPUs, and the remaining \(G_{\text{remaining}}\) GPUs will be allocated in subsequent steps.

\subsubsection{Island Distribution}

In the case where \(G_{\text{remaining}} = 0\), all GPUs have already been assigned to the \(N_{\text{actual}}\) islands at their minimum size, so we make no weighted adjustments.
When \(N_{\text{actual}} = 1\), there is only a single island so we add all \(G_{\text{remaining}}\) GPUs to that island.
Otherwise we can distribute the remaining GPUs using the \(S_{\text{skew}}\)value.

We first compute a weight $w_k$ for each island $k$, where $k$ ranges from $1$ to $N_{\text{actual}}$, computed as $w_k = k^{S_{\text{skew}}}$.
However, if $S_{\text{skew}}$ is numerically close to zero, each island is given an equal weight of $w_k = 1.0$.
These weights are then normalised so that we can derive proportions $p_k$, indicating each island's share of the remaining GPUs. The sum of the weights is computed as
\[
W_{\text{sum}} = \sum_{j=1}^{N_{\text{actual}}} w_j.
\]
%If $W_{\text{sum}} \leq 0$, which may occur due to numerical errors or degenerate inputs, all islands are assigned equal proportions, $p_k = \frac{1}{N_{\text{actual}}}$. Otherwise, 
Each proportion is determined by dividing the weight of an island by the weight sum:
\[
p_k = \frac{w_k}{W_{\text{sum}}}.
\]
Using these proportions, we compute the ideal (possibly fractional) allocation $s_k^{\text{ideal}}$ for each island:
\[
s_k^{\text{ideal}} = p_k \cdot G_{\text{remaining}}.
\]

\subsubsection{Integer Adjustment}

Because GPUs can only be allocated in integer units, we use the largest remainder method to convert the fractional allocations into integers.
First, each island $k$ is provisionally assigned a base number of GPUs:
\[
s_k^{\text{base\_int}} = \left\lfloor s_k^{\text{ideal}} \right\rfloor.
\]
The current island sizes (initialised to $I_{\text{min}}$) are then increased by these integer allocations.

%\todo{
%
%Next, the discrepancy between the total allocated GPUs and the intended $G_{\text{remaining}}$ is computed. This discrepancy,
%\[
%D = \text{round} \left( G_{\text{remaining}} - \sum_{j=1}^{N_{\text{actual}}} s_j^{\text{base\_int}} \right),
%\]
%represents the number of GPUs that are still unassigned due to rounding down the fractional values. To distribute the $D$ leftover GPUs, we calculate the fractional part for each island as
%\[
%f_k = s_k^{\text{ideal}} - s_k^{\text{base\_int}},
%\]
%and identify the $D$ islands with the largest fractional parts. Each of these selected islands receives one additional GPU.
%
%At the end of this process, each island has a final size given by
%\[
%S_k = I_{\text{min}} + s_k^{\text{base\_int}} + \delta_k,
%\]
%where $\delta_k \in \{0, 1\}$ depending on whether the island received an extra GPU during the discrepancy distribution. The total sum of all final island sizes equals $G_{\text{total}}$, satisfying the constraint that all available GPUs are used. The function returns this list of $N_{\text{actual}}$ island sizes as the result of the procedure.
%
%}


\clearpage
\subsubsection*{Search Parameters}

\begin{itemize}
    \item \( T \): the set of GPU types, e.g., \( T = \{\text{A100}, \text{V100}, \ldots\} \)
    \item \( n_t \): the target number of islands for GPU type \( t \in T \)
    \item \( s_t \): a vector value representing the skew or balance in GPU allocation across the \( n_t \) islands for GPU type \( t \)
    \item \( I_t \): the total available inventory of GPU type \( t \)
\end{itemize}

%Here, \( n_g \) is an integer representing how many islands to distribute GPUs across, and \( s_g \in [0, 1] \) controls how evenly the GPUs are distributed across those islands. For example, \( s_g = 0 \) could represent perfectly balanced islands (equal GPU count per island), while \( s_g = 1 \) could represent maximal skew (e.g., one island with most of the GPUs, others with minimal or none).

These two parameters (\( n_t \) and \( s_t \)) define the entire resource distribution strategy, allowing the Bayesian Optimisation process to search over high-level allocation patterns while implicitly satisfying inventory constraints by design.

\subsubsection*{Implementation}

\todo{
\begin{itemize}
\item Ax Implementation.
\item SAASBO not required.
\end{itemize}
}
